
%\documentclass{amsart}
\documentclass[12pt]{amsart}

\usepackage[margin=2cm]{geometry}
\usepackage{latexsym,amsfonts,amssymb,amsthm,amsmath,amscd,stmaryrd}
\usepackage{mathrsfs}
\usepackage[all]{xy}  
\usepackage{enumerate}
\usepackage{color,tikz}
\usepackage{multicol}
\usepackage{empheq}
\usepackage{comment}
\usepackage{enumitem}

\usepackage{tikz-cd}



%For sets
\newcommand{\set}[1]{\left\{ #1 \right\}}

%Letters
\newcommand{\KK}{\mathbb{K}}

\newcommand{\RR}{\mathbb{R}}
\newcommand{\NN}{\mathbb{N}}
\newcommand{\ZZ}{\mathbb{Z}}
\newcommand{\QQ}{\mathbb{Q}}
\newcommand{\FF}{\mathbb{F}}
\newcommand{\CC}{\mathbb{C}}
\newcommand{\m}{\mathfrak{m}}
\newcommand{\V}{\mathcal{V}}
\newcommand{\p}{\mathfrak{p}}

%operators
\DeclareMathOperator{\Ring}{\textbf{Ring}}
\DeclareMathOperator{\ev}{ev}

\newcommand{\syz}{\operatorname{syz}}
\newcommand{\Spec}{\operatorname{Spec}}
\newcommand{\Depth}{\operatorname{depth}}
\newcommand{\Hom}{\operatorname{Hom}}
\newcommand{\End}{\operatorname{End}}
\newcommand{\Ext}{\operatorname{Ext}}
\newcommand{\Tor}{\operatorname{Tor}}
\renewcommand{\ker}{\operatorname{ker}}
\newcommand{\coker}{\operatorname{coker}}
\newcommand{\Max}{\operatorname{max}}
\newcommand{\Ann}{\operatorname{Ann}}
\newcommand{\Ass}{\operatorname{Ass}}	
\newcommand{\Char}{\operatorname{char}}
\newcommand{\rk}{\operatorname{rk}}	
\newcommand{\height}{\operatorname{ht}}	
\newcommand{\ceil}[1]{\lceil {#1} \rceil}
\newcommand{\floor}[1]{\lfloor {#1} \rfloor}
\newcommand{\ann}{\operatorname{ann}}
\newcommand{\Min}{\operatorname{Min}}
\DeclareMathOperator{\im}{im}
\renewcommand{\dim}{\operatorname{dim}}
\newcommand{\grade}{\operatorname{grade}}
\newcommand{\reg}{\operatorname{reg}}
\newcommand{\depth}{\operatorname{depth}}
\DeclareMathOperator{\Supp}{Supp}
\newcommand{\id}{\operatorname{id}}
\newcommand{\Tr}{\operatorname{Tr}}
\renewcommand{\height}{\operatorname{ht}}
\newcommand{\Rmod}{\operatorname{R-mod}}
\newcommand{\pd}{\operatorname{pd}}


%leg/geqslant
\newcommand{\ls}{\leqslant}
\newcommand{\gs}{\geqslant}

%(co)homology
\renewcommand{\H}{\operatorname{H}}
%Cech
\newcommand{\Ch}{\textrm{\v C}}


%colors
\newcommand{\blue}[1]{\color {blue} #1}                   
\newcommand{\red}[1]{\color {red} #1}                   
\newcommand{\green}[1]{\color {green} #1}  

%primes


\begin{document}

	\thispagestyle{empty}\kern -2em
\rightline{\emph{\tiny (c) Juliette Bruce 2025}}
\kern -.5em
\rightline{\emph{\tiny licensed under a Creative Commons}}
\kern -.5em
\rightline{\emph{\tiny By-NC-SA 4.0 International License.}}
% \rightline{\includegraphics[width=.75in]{cc_notice}}

\bigskip
	
\begin{center}
	{\Large Worksheet 1.2: Quotients, Primes, Maximal Ideals}
\end{center}


\medskip

\noindent	 Let $R$ be any ring.  In this course, ``ring" means  {\it commutative ring with unity}.

\medskip

	\noindent	 \hrulefill 
	
\begin{enumerate}
	\item \textbf{Operations on Ideals.} Throughout this problem let $I$ and $J$ be ideals in a ring $R$.
	\begin{enumerate}[(a)]
		\item Prove or disprove that $I \cap J$ is an ideal. 
		\item Prove or disprove that $I \cup J$ is an ideal.
		\item Prove or disprove that the product of ideals, defined below, is an ideal:
		\[
		IJ \coloneqq \left\{ a_1b_1+\cdots+a_tb_t\;\; \big| \;\; a_i\in I\;\; b_i \in J, \; \; \text{for all $t$}\right\}.
		\]
		\item Prove or disprove that the sum of ideals, defined below, is an ideal:
		\[
		I + J \coloneqq \left\{ a + b \;\; \big| \;\; a \in I\;\; b \in J\right\}.
		\]
		\item Prove or disprove that the radical of an ideal, defined below, is an ideal:
		\[
		\sqrt{I} \coloneqq \left\{ a \in R \;\; \big| \;\; a^{k} \in I \text{ for some $k\in \ZZ_{\geq1}$} \right\}.
		\]
	\end{enumerate}
	
	\ \\
	
	\item Let $\KK$ be a field. The rings $\ZZ$, $\KK[x]$, and $\KK[[t]]$ are all principal ideal domains (PIDs) meaning domains where every ideal is generated by a single element. If $I= \langle a \rangle$ and $J = \langle b \rangle$ give, with proof, a careful concrete description of generators for $IJ$, $I\cap J$, $I+J$, and $\sqrt{I}$ in each of these rings.

\end{enumerate}
\ \\

\noindent \textbf{Universal Properties.} Throughout this course we will see numerous constructions and operations (e.g., quotients, tensor products, localizations, completions, etc.) each of these will have an explicit construction where we describe the objects precisely as sets. However, there is an alternative perspective, which is just as useful, describing these objects via their \emph{universal properties}, which roughly asks what makes each object special amongst all other such objects. The standard framework for a universal property is, ``Object X is unique up to isomorphism among all objects having property P.''
\begin{enumerate}[start=3] 

	\item \textbf{Universal Property of Quotients.} Let $R$ be a ring and $I \subseteq R$ an ideal. Let $\pi\colon R \to R/I$ be the quotient map. Suppose $\varphi \colon R \to S$ is a ring homomorphism such that $I \subseteq \ker(\varphi)$. Define a function
    \[
    \widetilde{\varphi} \colon R/I \to S
    \qquad \text{by} \qquad
    \widetilde{\varphi}(r+I)=\varphi(r).
    \]
    \begin{enumerate}[(a)]
    	\item Prove that $\widetilde{\varphi}$ is well-defined.

    	\item Prove that $\widetilde{\varphi}$ is a ring homomorphism.

    	\item Prove that $\widetilde{\varphi} \circ \pi = \varphi$.

    	\item Prove that $\widetilde{\varphi}$ is the unique ring homomorphism $R/I \to S$ with the property that
    	\[
  	  \widetilde{\varphi}\circ \pi = \varphi.
    	\]
    	\item Explain in words what this says about the quotient ring $R/I$.
    \end{enumerate}

\item In this problem, we wish to consider the evaluation map: 
    \[
    \ev_i \colon \ZZ[x] \to \CC,
    \qquad
    f(x)\mapsto f(i)
    \]
    \begin{enumerate}[(a)]
    \item Prove that $\ev_{i}$ is a ring homomorphism.

    \item Show that $\langle x^2+1 \rangle \subset \ker(\ev_i)$.

    \item Use Problem 3 to produce a ring homomorphism
    \[
    \widetilde{\ev_i}\colon \ZZ[x]/\langle x^2+\rangle \to \CC.
    \]

    \item What does $\widetilde{\ev_i}$ do to the class of $x$ in $\ZZ[x]/\langle x^2+1\rangle$?
    
    \item Prove that $\ZZ[x]/\langle x^2+1\rangle$ is isomorphic to $\ZZ[i]$. Is $\langle x^2 +1\rangle$ prime and/or maximal?
\end{enumerate}
\ \\

\item \textbf{Universal Property of Polynomial Rings.} The goal of this exercise is to prove the following characterization of polynomial rings. Let $\varphi\colon R \to S$ be a ring homomorphism. For any element $s \in S$ there is a unique ring homomorphism $\psi\colon R[x] \to S$ such that $\psi(x) = s$ and the following diagram commutes:
\[
\begin{tikzcd}[row sep = 3.5em, column sep = 3.5em]
	R \arrow[r, "\varphi"] \arrow[d, hook] & S  \\
	R[x] \arrow[ur, dashed, swap,"\psi"]
\end{tikzcd}.
\]
\begin{enumerate}[(a)]
    \item Prove there is at most one ring homomorphism $\psi:R[x] \to S$ with the above properties. 
    \item Construct such a homomorphism explicitly by setting
    \[
    \psi\!\left(\sum_{i=0}^n a_i x^i\right)=\sum_{i=0}^n \varphi(a_i)s^i.
    \]
    \item Conclude that to give a ring homomorphism $R[x]\to S$ is equivalent to giving:
    \begin{itemize}
        \item a ring homomorphism $R\to S$, and
        \item an element of $S$.
    \end{itemize}
\end{enumerate}
\ \\

\item \textbf{Evaluation maps.} Let $R$ be a ring and let $a\in R$. Define the evaluation map
    \[
    \ev_a \colon R[x]\to R,
    \qquad
    f(x)\mapsto f(a).
    \]
    \begin{enumerate}[(a)]
    \item Prove that $\ev_a$ is a surjective ring homomorphism.

    \item Prove that $\ker(\ev_a) = \langle x-a\rangle$.

    \item Use Problem 1 to show that $\ev_a$ factors through a ring homomorphism
    \[
    R[x]/(x-a)\to R.
    \]
    
    \item Explain what this means if $R$ is a field or $R$ is an integral domain.
    
    \item Give an example of a ring $R$ and prime/maximal ideal $I \subset R[x]$ such that $I \neq \ker(\ev_a)$ for all $a \in R$. 
\end{enumerate}

\end{enumerate}
\end{document}





